\clearpage\pdfbookmark[0]{技术附录}{appendix}\section*{技术附录}
\setcounter{table}{0}\renewcommand{\thetable}{A\arabic{table}}\setcounter{figure}{0}\renewcommand{\thefigure}{A\arabic{figure}}
\subsection*{A.1 数据构造与模型选择}
主面板以客户为单位划分开发与评价数据。一般历史开发池包含128,069个快照、49,371名客户；提额场景开发池包含22,524个状态快照，覆盖事件前一期、提额期和相邻期。拟合与选择采用客户级三折划分，校准客户单独留出。最终拟合使用107,880个结局已知快照，涉及37,974名客户、2,304个阳性窗口。一般历史与提额场景各占一半总训练权重，各来源内再按客户有效快照数的倒数加权。校准使用1,112个提额期已知结果，其中30例事件。

244维表示由179项历史特征和65项现金流特征构成，覆盖余额、支用、支付、最低应还及其变化和窗口统计。414维表示进一步加入170项滞后、比例与多窗口特征。参考视图从原始历史按统一额度尺度重算；缺失处理仅在训练数据上拟合，输入截止于评分时点。各模型预测按客户、账户和观察期逐行对齐。

\begin{table}[H]\centering\caption{强基线配置与拟合范围}\small
\begin{tabularx}{\linewidth}{@{}p{23mm}YY@{}}\toprule
模型族 & 主要参数 & 表示与拟合次数\\\midrule
CatBoost & 600轮；学习率0.025；深度5；L2=15；列采样0.9 & 实际／参考244维，各三折＋最终拟合\\
LightGBM & 600轮；学习率0.025；31叶；深度6；叶节点最少100；L2=20、L1=0.5 & 实际／参考244维及参考414维，各三折＋最终拟合\\
XGBoost & 600轮；学习率0.025；深度4；子叶权重8；L2=20、L1=0.5 & 实际／参考244维，各三折＋最终拟合\\\bottomrule
\end{tabularx}\tabnote{七配置、28次拟合，随机种子20261003。LightGBM与XGBoost行／列采样0.9；CatBoost行采样采用实际默认设置。}\end{table}

开发比较覆盖单模型、等权集成和异构混合，配合静态、指数及限幅更新，共形成126个候选。筛选同时约束风险识别与工作容量：两期平均AUC不低于S减0.005，AP不低于S的98\%，Brier和Logloss不高于S的102\%；10\%容量下两期覆盖之和不低于S，任一期最多少2例；六档容量平均召回不低于S的98\%。合格候选再依较差期CVaR95、AP和覆盖排序。S为开发平均AP最高的实际尺度静态集成；R为244维目标加权管线，T为既有尾部优化方案，FX、E分别为参考混合、参考集成的条件指数更新。

本文固定选定的244维XGBoost与414维LightGBM等权底模及其校准，在开发交叉拟合预测上比较容量保护方式。主作业容量预定为10\%，另考察同时保护六档边界的方案。两阶段分别包含6,411、6,308条有效更新；主版本平均AUC为0.8004、AP为0.1672，覆盖173个事件窗口，较差期CVaR95为1.612个百分点。选择结果见A.10。

\subsection*{A.2 终点、评价单位与区间}
主终点为评分后三个相对月份内已观察DPD$>0$。观察到阳性即记为事件；随访不足且尚未观察到事件者记为未知。复核容量按全部客户计算，结果未知客户同样占用名额。两期分别有99、88个事件窗口；重叠窗口可能涉及同一客户，独立事件客户另行去重。三个月内已知DPD30事件为0；6／12个月延长随访结果见A.5。

历史评价队列曾用于研究迭代。本次比较固定底模与逐行预测，区间刻画该队列的客户抽样变异；真实日历外推与新机构一次性评价采用独立协议。尾部比较按客户重采样2,000次，三客群、两期、七对手共42项比较共享权重，以最大标准化中心偏差构造同时95\%区间，临界值为3.8839。固定名单覆盖差和原106条更新子组分别报告逐项区间。上述条件区间未纳入开发选择的不确定性。

\clearpage\subsection*{A.3 完整风险指标与校准表现}
\begin{table}[H]\centering\caption{完整风险指标与概率校准诊断}
\fontsize{9bp}{12bp}\selectfont
\begin{tabularx}{\linewidth}{@{}Yrrrrrr@{}}\toprule
阶段 / 方法 & KS & Gini & Brier & Logloss & 截距 & 斜率\\\midrule
提额月 / 目标加权 R & 0.5059 & 0.6246 & 0.03231 & 0.13138 & -0.167 & 0.847\\
提额月 / 风险集成 S & 0.5164 & 0.6385 & 0.03191 & 0.12978 & -0.092 & 0.875\\
提额月 / 尾部优化 T & 0.5157 & 0.6329 & 0.03177 & 0.12965 & -0.038 & 0.892\\
提额月 / 参考混合 FX & 0.5080 & 0.6342 & 0.03188 & 0.12971 & -0.026 & 0.899\\
提额月 / 参考集成 E & 0.5264 & 0.6376 & 0.03199 & 0.13019 & -0.069 & 0.884\\
提额月 / 同底模静态 & 0.5031 & 0.6348 & 0.03186 & 0.12963 & -0.006 & 0.904\\
提额月 / 原1pp更新 & 0.5023 & 0.6336 & 0.03189 & 0.12979 & 0.027 & 0.911\\
提额月 / 本文方法 & 0.5023 & 0.6342 & 0.03186 & 0.12968 & 0.042 & 0.916\\
提额月 / 六容量保护 & 0.5023 & 0.6344 & 0.03186 & 0.12963 & 0.040 & 0.916\\
相邻月 / 目标加权 R & 0.4665 & 0.5351 & 0.03140 & 0.13514 & -0.038 & 0.849\\
相邻月 / 风险集成 S & 0.4596 & 0.5492 & 0.03136 & 0.13628 & 0.115 & 0.862\\
相邻月 / 尾部优化 T & 0.4354 & 0.5427 & 0.03108 & 0.13342 & -0.159 & 0.834\\
相邻月 / 参考混合 FX & 0.4437 & 0.5439 & 0.03104 & 0.13268 & -0.117 & 0.851\\
相邻月 / 参考集成 E & 0.4325 & 0.5473 & 0.03118 & 0.13336 & -0.173 & 0.836\\
相邻月 / 同底模静态 & 0.4437 & 0.5405 & 0.03107 & 0.13294 & -0.096 & 0.854\\
相邻月 / 原1pp更新 & 0.4425 & 0.5400 & 0.03105 & 0.13289 & -0.096 & 0.855\\
相邻月 / 本文方法 & 0.4425 & 0.5403 & 0.03102 & 0.13279 & -0.076 & 0.860\\
相邻月 / 六容量保护 & 0.4429 & 0.5405 & 0.03103 & 0.13284 & -0.077 & 0.859\\
\bottomrule\end{tabularx}
\tabnote{有效样本分别为2,841、2,623人。截距与斜率在已知结局上作描述性拟合，理想值为0、1，诊断结果不回写评分。}
\end{table}
\fig{v18_reliability}{最终评分的概率校准。各方法分别划分十个等人数组；本文误差线为组内事件率的Wilson 95\%区间。}
两期AUC、AP、KS、Gini、概率损失与校准均按最终$q$计算。校准截距、斜率由评价样本拟合$\operatorname{logit}\Pr(Y=1)=a+b\operatorname{logit}(q)$得到，理想值为$(0,1)$。拟合仅用于诊断，预测保持冻结；等人数组区间刻画组内事件率的不确定性。

\clearpage\subsection*{A.4 容量范围与连续风险覆盖}
\begin{table}[H]\centering\caption{两期六档固定容量的已观察事件覆盖}
\fontsize{9bp}{12bp}\selectfont
\begin{tabularx}{\linewidth}{@{}Yrrrrrr@{}}\toprule
阶段 / 方法 & 5\% & 8\% & 10\% & 12\% & 15\% & 20\%\\\midrule
提额月 / 目标加权 R & 27 & 40 & 48 & 54 & 60 & 64\\
提额月 / 风险集成 S & 26 & 42 & 44 & 52 & 56 & 65\\
提额月 / 尾部优化 T & 28 & 37 & 43 & 50 & 55 & 64\\
提额月 / 参考混合 FX & 28 & 37 & 45 & 51 & 55 & 63\\
提额月 / 参考集成 E & 26 & 40 & 44 & 51 & 56 & 65\\
提额月 / 同底模静态 & 27 & 39 & 46 & 52 & 55 & 63\\
提额月 / 原1pp更新 & 28 & 38 & 44 & 51 & 55 & 63\\
提额月 / 本文方法 & 28 & 38 & 46 & 51 & 55 & 63\\
提额月 / 六容量保护 & 27 & 39 & 46 & 52 & 55 & 63\\
相邻月 / 目标加权 R & 19 & 25 & 32 & 37 & 42 & 49\\
相邻月 / 风险集成 S & 15 & 28 & 34 & 35 & 41 & 48\\
相邻月 / 尾部优化 T & 20 & 27 & 32 & 32 & 39 & 48\\
相邻月 / 参考混合 FX & 17 & 28 & 32 & 35 & 42 & 51\\
相邻月 / 参考集成 E & 16 & 26 & 33 & 35 & 40 & 50\\
相邻月 / 同底模静态 & 16 & 29 & 33 & 35 & 42 & 51\\
相邻月 / 原1pp更新 & 17 & 28 & 32 & 35 & 42 & 51\\
相邻月 / 本文方法 & 17 & 29 & 33 & 35 & 42 & 51\\
相邻月 / 六容量保护 & 16 & 29 & 33 & 35 & 42 & 51\\
\bottomrule\end{tabularx}
\tabnote{各容量名额向上取整。主版本按10\%边界生成一次评分，再评价六档容量；六容量版本一次保护全部六个边界。}
\end{table}
\fig{v21_capacity_range}{不同复核容量下的风险覆盖。上部比较两期六档容量覆盖的已观察事件数；下部给出本文与S、T的覆盖差。主版本保护10\%边界，其余容量检验同一评分对工作量变化的适应性。}

\clearpage\subsection*{A.4（续）近稳口径与客户接续}
\begin{table}[H]\centering\caption{账务近稳客户的评分波动}
\fontsize{9bp}{12bp}\selectfont
\begin{tabularx}{\linewidth}{@{}Yrrrr@{}}\toprule
 & \multicolumn{2}{c}{提额月：582对} & \multicolumn{2}{c}{相邻月：439对}\\方法 & P95 & CVaR95 & P95 & CVaR95\\\midrule
目标加权 R & 7.466 & 14.014 & 1.418 & 4.389\\
风险集成 S & 6.107 & 10.994 & 1.594 & 4.299\\
尾部优化 T & 1.436 & 3.634 & 0.470 & 1.535\\
参考混合 FX & 1.595 & 3.978 & 0.608 & 1.890\\
参考集成 E & 1.518 & 4.049 & 0.494 & 1.806\\
同底模静态 & 2.994 & 7.378 & 1.007 & 3.386\\
原1pp更新 & 1.000 & 1.332 & 1.000 & 1.465\\
\textbf{本文方法} & 1.000 & 2.184 & 1.000 & 1.506\\
六容量保护 & 1.261 & 3.272 & 1.000 & 1.967\\
\bottomrule\end{tabularx}
\tabnote{单位为概率百分点。近稳评价使用余额、支付、支用、最低应还四个字段；更新门控还包含三项应收金额。}
\end{table}
\begin{table}[H]\centering\caption{不同四字段近稳门槛下的表现}
\fontsize{9bp}{12bp}\selectfont
\begin{tabularx}{\linewidth}{@{}Yrrrrr@{}}\toprule
阶段 / 阈值 & 对数 & 均值 & P95 & CVaR95 & 超过2 pp\%\\\midrule
提额月 / 1\% & 137 & 0.394 & 1.000 & 3.974 & 1.46\\
提额月 / 5\% & 582 & 0.368 & 1.000 & 2.184 & 0.86\\
提额月 / 10\% & 899 & 0.416 & 1.264 & 3.066 & 2.56\\
相邻月 / 1\% & 134 & 0.164 & 0.844 & 0.994 & 0.00\\
相邻月 / 5\% & 439 & 0.229 & 1.000 & 1.506 & 0.23\\
相邻月 / 10\% & 662 & 0.325 & 1.000 & 2.805 & 1.96\\
\bottomrule\end{tabularx}
\tabnote{幅度均为概率百分点；仅改变评价分组，最终评分及更新规则保持冻结。}
\end{table}
四字段近稳组依据余额、总支付、总支用和最低应还划分；七字段更新条件另含应收本金、应收总额与应收金额。更新门控同时要求字段完整、同账户同周期、月份连续、两期DPD为0及至少六期历史。本文两期10\%名单与同底模静态名单逐人一致；六容量版本在全部六档容量下均保持成员一致。

\clearpage\subsection*{A.4（续）固定名额的去重结果与区间}
\begin{table}[H]\centering\caption{两期559个名额的客户接续与事件覆盖}
\fontsize{9bp}{12bp}\selectfont
\begin{tabularx}{\linewidth}{@{}Yrrrr@{}}\toprule
方法 & 独立客户 & 独立事件客户 & 连续入选 & 相邻月新增\\\midrule
目标加权 R & 356 & 61 & 203 & 65\\
风险集成 S & 347 & 60 & 212 & 56\\
尾部优化 T & 345 & 55 & 214 & 54\\
参考混合 FX & 329 & 60 & 230 & 38\\
参考集成 E & 337 & 60 & 222 & 46\\
同底模静态 & 331 & 61 & 228 & 40\\
原1pp更新 & 328 & 59 & 231 & 37\\
\textbf{本文方法} & 331 & 61 & 228 & 40\\
六容量保护 & 331 & 61 & 228 & 40\\
\bottomrule\end{tabularx}
\tabnote{连续入选客户需要核查新账务；原材料、处理意见及责任人随原案接续。}
\end{table}
\begin{table}[H]\centering\caption{两期事件窗口覆盖差的配对区间}
\fontsize{9bp}{12bp}\selectfont
\begin{tabularx}{\linewidth}{@{}Yrr@{}}\toprule
对手 & 本文减对手 & 逐项95\%区间\\\midrule
目标加权 R & -1 & [-11, 9]\\
风险集成 S & 1 & [-8, 10]\\
尾部优化 T & 4 & [-2, 11]\\
参考混合 FX & 2 & [0, 5]\\
参考集成 E & 2 & [-6, 9]\\
同底模静态 & 0 & [0, 0]\\
原1pp更新 & 3 & [0, 7]\\
\bottomrule\end{tabularx}
\tabnote{共559个名额，本文覆盖79个事件窗口。名单固定，按客户贡献重采样2,000次；同一客户在两期共享权重。窗口数可重复，独立事件客户另表列示。}
\end{table}
容量保护恢复原更新遗漏的三个事件窗口：提额月两例、相邻月一例。两期独立事件客户由59增至61；559个名额涉及331名独立客户，其中228名连续入选。事件覆盖和独立接触人数反映名单安排效果，实际工时由现场办理记录计量。

\clearpage\subsection*{A.5 边界更新的风险表现与严重风险随访}
\begin{table}[H]\centering\caption{原106次受限更新的条件评价}
\fontsize{9bp}{12bp}\selectfont
\begin{tabularx}{\linewidth}{@{}Yrrrrrrr@{}}\toprule
方法 & 均值 & AUC & AP & Brier & Logloss & 偏移截距 & 斜率\\\midrule
同底模静态 & 9.90\% & 0.7088 & 0.2354 & 0.09810 & 0.33276 & 0.179 & 0.747\\
原1pp更新 & 9.12\% & 0.7079 & 0.2531 & 0.09832 & 0.33578 & 0.277 & 0.786\\
\textbf{本文方法} & 8.74\% & 0.7384 & 0.2843 & 0.09704 & 0.33034 & 0.328 & 0.894\\
\bottomrule\end{tabularx}
\tabnote{106次更新、98名客户；105次结果可判、12个事件。偏移截距将预测logit作为offset，斜率来自截距与斜率共同拟合。}
\end{table}
\begin{table}[H]\centering\caption{固定受限子组的配对变化}
\fontsize{9bp}{12bp}\selectfont
\begin{tabularx}{\linewidth}{@{}Yrr@{}}\toprule
对手 / 指标 & 本文减对手 & 逐项95\%区间\\\midrule
同底模静态 / brier & -0.00107 & [-0.00576, 0.00409]\\
同底模静态 / logloss & -0.00242 & [-0.01975, 0.01634]\\
同底模静态 / auc & 0.02957 & [-0.02578, 0.08806]\\
同底模静态 / ap & 0.04884 & [-0.01564, 0.16210]\\
原1pp更新 / brier & -0.00128 & [-0.00362, -0.00002]\\
原1pp更新 / logloss & -0.00544 & [-0.01453, -0.00019]\\
原1pp更新 / auc & 0.03047 & [0.00000, 0.07264]\\
原1pp更新 / ap & 0.03118 & [0.00000, 0.10919]\\
\bottomrule\end{tabularx}
\tabnote{以原方案的106条记录固定子组，客户成组重采样2,000次；这些区间属于条件比较，不与尾部42项共同校正。}
\end{table}
对原方案实际修正至边界的106次更新单独评价，涉及98名客户，其中105次结果可判、12个事件。容量保护降低了该组Brier和Logloss。该组平均预测为8.74\%，事件率为11.43\%；8次阳性向上受限记录的预测仍低于静态值，其中2次得到部分恢复。名单保持、概率校准与条件风险响应由此分别检验。

从完整批次看，容量保护恢复了三个事件窗口的入选资格，其名次由319、306、271分别变为286、288、267，对应提额月291个名额和相邻月268个名额。第三例本身未触发限幅，而是因其他客户分数变化被挤出。整批边界处理同时解决客户自身受限与名单内相互挤出两类问题。

\clearpage\subsection*{A.5（续）极端更新与延长终点}
5,584次更新中，容量保护对14次近稳候选作必要放宽，提额月12次、相邻月2次，涉及14名客户。最大额外放宽为18.958个百分点：前次风险26.016\%，基础预测4.077\%，近稳候选25.016\%，最终6.058\%。基础模型将该客户排在名单之外，容量保护据此将读数移回未入选边界，总跨期变化为19.958个百分点。极端记录与逐户原因共同补充P95及平均幅度，明确稳定更新的实际边界。

\begin{table}[H]\centering\caption{延长随访的DPD30+结果}
\fontsize{9bp}{12bp}\selectfont
\begin{tabularx}{\linewidth}{@{}Yrrrrr@{}}\toprule
期限 / 方法 & 结局已知 & 事件窗口 & 事件客户 & 入选窗口 & 入选客户\\\midrule
6个月 / 同底模静态 & 4417 & 5 & 3 & 2 & 1\\
6个月 / 原1pp更新 & 4417 & 5 & 3 & 2 & 1\\
6个月 / 本文方法 & 4417 & 5 & 3 & 2 & 1\\
12个月 / 同底模静态 & 3138 & 16 & 9 & 3 & 2\\
12个月 / 原1pp更新 & 3138 & 16 & 9 & 3 & 2\\
12个月 / 本文方法 & 3138 & 16 & 9 & 3 & 2\\
\bottomrule\end{tabularx}
\tabnote{沿用原559个名额和冻结评分；遇到DPD30+即记事件，未观察到事件且不足全窗口者记未知。延长终点不重新解释为三个月PD。}
\end{table}
6个月随访观察到5个DPD30+窗口、3名客户；12个月观察到16个窗口、9名客户。三种同底模版本在原容量内分别覆盖2／3个窗口、1／2名客户。严重事件样本较少，结果按窗口数和独立人数呈现；三个月评分仍对应早期风险终点。外部数据的真实核销与逾期结果见A.7。

\subsection*{新机构验证的登记与随访}
新机构验证先按客户排除历次研究记录，登记模型、特征、容量规则和预测，待结局成熟后统一评价。数据合同纳入真实日历时点、动作原因、审批发起方与资料可用日期。程序记录首次读取结局的状态，并校验重复读取、日期倒置和随访未成熟记录。历史研究与机构一次性验证各自形成独立结果。

同源数据排除历次客户后，剩余4,998名客户、16,734条记录，包含380次连续正额度增加；沿用动作前至少六期历史、前后未逾期等条件，合格新事件为0。公开数据的23个原始字段未记录动作原因和真实日历日期，因此识别的是观察提额；经营性提额在机构验证中依据审批记录和策略凭证认定。

\clearpage\subsection*{A.6 参考尺度与容量保护的数学性质}
\textbf{命题1（固定批次下的纯尺度不变性）。}固定所有客户的非额度历史、参考周期、前次状态、模型、缺失模式及容量。参考构建器和门控均不读取当期额度作为变量，基础预测、近稳候选与排序边界随之不变。故仅改变当期额度不会改变$q$。客户账务更新、参考周期重建或批次成员变化时重新计算。

\textbf{命题2（预定容量的成员保持）。}以基础预测降序、客户行键升序定义确定次序。无边界并列时，取第$k$与$k+1$名预测的中点$\tau$，并令$\tau^+$为其向上相邻可表示数。入选客户的允许区间为$J_i=[\tau^+,1]$，其余为$J_i=[0,\tau]$。令$q_i=\Pi_{J_i}(z_i)$，则所有入选分数严格大于未入选分数，因而$\operatorname{Top}_k(q)=\operatorname{Top}_k(p)$。并列时按同一行键施加至多$2\times10^{-10}$的有序概率微扰；边界0与1采用向内部留出数值空间的规则。

\textbf{命题3（给定名单区间内的最小修正）。}区间$J_i$固定时，$\sum_i(u_i-z_i)^2$在$\prod_iJ_i$上的唯一最小解为逐坐标投影$q_i=\Pi_{J_i}(z_i)$。原值可行时不变，否则取最近端点。该最小性针对指定分界区间；名单之外的完整排序无需保持。若同时保护六个容量边界，则按基础排序将客户划入相应层区间，逐层应用同一投影。

近稳候选$z$由$q_i^-$的$\pm1$个百分点区间生成；最终$q$优先满足名单边界，必要时放宽。一个客户的输出因而依赖本次完整批次、名额及前次状态。接口用批次、模型、参考周期和容量共同标识版本；修改成员或容量后整批重算，批次内状态原子写入。8,491个历史状态重放与冻结结果最大差为$1.67\times10^{-16}$。

\subsection*{参考周期与路径解释}
令$z(H;l)=\operatorname{logit}p(H;l)$。在$H_0$处切换额度，以旧状态偏移消除当时跳变后，下一期仍存在交互项
\[
\Omega=z(H_1;l_1)-z(H_0;l_1)-z(H_1;l_0)+z(H_0;l_0).
\]
完整乘积域上，所有四角的$\Omega=0$当且仅当$z(H;l)=a(H)+b(l)$：固定$(H_*,l_*)$，令$a(H)=z(H;l_*)$，$b(l)=z(H_*;l)-z(H_*;l_*)$代入即可。不完整二部支持图须检查每个连通分量全部闭环的交替和；仅检查已存在四角不足以覆盖长环。

\subsection*{路径界与变化分解}
百分制读数$V(H;l)=100[1-p(H;l)]$在旧偏移与规范参考两条路径间满足
\[
|V_{\rm stale}-V_{\rm canonical}|\leq100\tanh(|\Omega|/4)\leq25|\Omega|.
\]
对$|\sigma(x+\Omega)-\sigma(x)|$最大化，两端关于0对称时取得等号。以当期状态重新计算桥接偏移$z(H_t;l_e)-z(H_t;l')$可恢复规范参考；产品采用周期内固定参考、周期变更显式初始化，并保留旧周期读数。

令$v_{ij}=100[1-p(H_i;l_j)]$，二元Shapley分量为
\[
\phi_H=\tfrac12[(v_{10}-v_{00})+(v_{11}-v_{01})],\qquad
\phi_l=\tfrac12[(v_{01}-v_{00})+(v_{11}-v_{10})].
\]
两者严格相加为$v_{11}-v_{00}$，表示两个更新次序的平均归属，解释对象为模型读数。四角组合须通过输入范围检查。周期内令$c_t=q_t-p_t$，实际显示变化满足
\[
Q_t-Q_{t-1}=100(p_{t-1}-p_t)+100(c_{t-1}-c_t).
\]
工作台将基础信息变化与更新修正分别显示；更新修正进一步保存近稳候选与容量保护两步。信息变化包含账务、窗口平移和缺失状态变化。实际敞口始终取当期额度与余额，参考额度不参与真实超限判断。
\clearpage\subsection*{A.7 参考周期机制与跨来源风险验证}
参考周期研究使用三族参考组件的2,677对同账户四角结果，检验非线性模型的路径交互。固定旧偏移时，下一期路径差P95为3.013分、最大26.676分；448人的更新方向相反，156人的路径差至少1分。该实验评价参考周期机制；正文统一评分的风险与尾部结果另由其冻结混合模型计算。

\fig{v18_reference_epoch}{参考组件的路径交互。a比较固定参考与固定旧偏移下的同客信息更新；b展示两条路径差的分位数。陶土色表示更新方向不同的448对记录。}

独立额度面板包含38,531名客户、每客14个月，形成469,836条相邻正额度转移、75,597次提额。使用率共变比例按客户簇重采样3,000次，另在留出客户和每客首次事件中计算。发布者内部评级用于描述研究范围，风险验证采用真实还款与核销标签。公开样本用于研究及界面演示，机构使用范围依各来源授权合同确定。

拍拍贷成熟窗口保留73,696个成熟可判定首标，其中16,446例事件。训练/验证/测试分别为13,722/17,681/10,347笔，对应事件3,452/4,068/2,099例。训练放款截至2015年末，窗口最晚结束2016年3月；验证放款2016年4至6月，窗口最晚结束9月；测试放款10月1日至11月22日，窗口截至2017年2月22日。源表在统一回溯时点发布，时间切分采用逻辑成熟窗口；首标及零历史成功借款作为客户隔离代理。

标签依据到期日、实际还款日、状态及正金额还款义务共同构造；未到期、零义务、计划不完整及日期冲突保留未知。正金额义务条件排除了训练期4笔零应还却标迟延记录，以及209笔窗口内无正金额应还记录；验证和测试标签未变。按统一标签规则与固定配置重新计算全部模型。

\begin{table}[H]\centering\caption{拍拍贷：全部固定模型的测试表现}
\fontsize{9.1bp}{12.5bp}\selectfont
\begin{tabularx}{\linewidth}{@{}Yrrrrrr@{}}
\toprule
模型 & AUC & AP & KS & Logloss & Brier & 10\%覆盖 \\
\midrule
Logistic & 0.5483 & 0.2266 & 0.0760 & 0.5036 & 0.1615 & 259 \\
CatBoost & 0.5566 & 0.2423 & 0.0766 & 0.5006 & 0.1605 & 282 \\
LightGBM & 0.5429 & 0.2325 & 0.0561 & 0.5056 & 0.1621 & 267 \\
XGBoost & 0.5535 & 0.2371 & 0.0721 & 0.5013 & 0.1609 & 274 \\
三族等权集成 & 0.5524 & 0.2388 & 0.0690 & 0.5016 & 0.1609 & 274 \\
\bottomrule\end{tabularx}
\tabnote{两个阶段各拟合4个模型：开发选择及训练+验证最终拟合。等权集成不另算训练。逐行标签条件及模型配置见复算材料。}
\end{table}

\clearpage
Prosper按LoanKey去重，排除冲突的ProsperScore字段，每个MemberKey保留最早合格贷款。征信日期须不晚于放款日，输入限于放款时已知信息。训练截至2009年11月，验证为2011年1至11月，测试自2013年1月起且观察月龄至少12，集合间隔至少13个月。训练、验证、测试分别有27,161人／3,739事件、6,363人／558事件和2,438人／77事件，借款人跨集合互斥。终点为前12还款周期核销。开发选定三族集成的测试AUC为0.7159、AP为0.0667，244人名单覆盖19例；Logistic、CatBoost各覆盖18例，LightGBM、XGBoost各19例。

新颜作者基准含66,533名训练客户、6,291名开发客户和12,412名日期后移留出客户，使用42项特征，各集合客户互斥。CatBoost、LightGBM、XGBoost各比较300、600树，缺失处理在训练集拟合。留出集中有2,005例还款异常；1,242个名额下，容量规则选定的Cat300与AP规则选定的Cat600均覆盖657例，等权集成覆盖660例。回连原始还款表后，部分开发标签的结果日期延伸至留出期，工程特征可用时点也未完整记录。该来源据此定位为作者预处理基准，严格成熟窗口评价采用上页拍拍贷协议。

\fig{v17_external_benchmarks}{跨来源风险识别与固定容量覆盖。新颜采用作者预处理基准，Prosper采用核销终点；各面板在对应事件率、时间窗口和容量下比较固定候选。}

\clearpage\subsection*{A.8 全量客户接续与产品接口}
工作台承载2,907名客户、17,441条当期可见月度记录，缺失月份保留为空，操作层不读取未来结局。相邻月2,677名客户构成完整容量保护批次，其中268人进入风险复核；189名较近观察期客户保留冻结评分，40人进入逾期核对，1人等待补账。全量客户分入五个互斥队列：名单复核268人、差异核对6人、持续观察2,592人、逾期核对40人、账务补齐1人。

189名补齐客户依据真实历史重算414项特征，输入截至观察期，沿用既定底模、校准和状态规则。已有冻结状态重放的最大差为$6.94\times10^{-17}$，独立重算最大概率差为$1.16\times10^{-16}$。任务记录与成熟评价队列分别保存；纳入新完整批次时，按该批次名额重新评分。

\begin{table}[H]\centering\caption{动作证据与参考周期的办理规则}\small
\begin{tabularx}{\linewidth}{@{}p{28mm}YY@{}}\toprule
情形 & 需核对的凭证 & 下一步\\\midrule
银行经营性调额 & 银行发起、资质未变、原始策略／审批凭证，且评分时已可用 & 沿用周期，保存依据与复审人\\
重新资质审查 & 新收入、资产或资质资料及其生效时间 & 交策略岗重建参考周期\\
账户或模型变更 & 账户关系、模型版本与切换批次 & 初始化新周期，保留前后记录\\
原因未知或凭证缺失 & 发起方、依据编号、证据可用日期 & 补件接续，待异岗复审\\\bottomrule
\end{tabularx}\end{table}

动作表单记录发起方、资质更新状态、凭证编号与证据可用观察月。沿用参考周期须同时满足“银行发起＋资质未变＋有效来源＋评分时已可用”，并由经办、复审两岗确认。原因未知的案件进入补件流程，待凭证齐备后认定动作类型。

评分适配器读取当期可见历史、参考周期、前次状态和整批容量，输出$p,z,q$、名单边界、修正原因与版本。案件服务按客户、账户、周期连接记录，通过版本检查处理并发、重复和乱序提交。Agent调用账务检索、比率计算、评分解释与材料检查工具，数值和状态统一取自服务结果；自然语言查询支持接入机构模型。来源引用、人工意见及处理回执随案保存，供下一岗位接续。

六个主页面、容量边界入口、268人名单次序、动作表单读写及Agent证据回连均已在软件执行环境完成验证。主文六页面集成图取自实际运行页面。机构接入阶段进一步联调统一身份、岗位权限、外部工单与异常恢复，并以真实作业结果验收。

\clearpage\subsection*{A.9 业务客群与尾部优势范围}
两类业务客群均按评分时可见账务划分：相邻两期账户均为Active、DPD为0，支付、最低应还及笔数字段完整且非负。\textbf{还款与交易状态组}将支付分为无最低应还、未覆盖最低应还、已覆盖最低应还三类，笔数分为0、1、2--5、6及以上四档；两期分类分别相同。\textbf{支付与交易结构组}要求最低应还均为正，支付／最低应还差不超过0.25，取现占支用份额差不超过0.10，笔数满足$2|n_1-n_0|/(1+n_0+n_1)\leq0.25$；取现与支用同为0时份额记0。

\begin{table}[H]\centering\caption{业务分组与七字段更新条件的重合}
\fontsize{9bp}{12bp}\selectfont
\begin{tabularx}{\linewidth}{@{}Yrrrrrr@{}}\toprule
 & \multicolumn{3}{c}{提额月} & \multicolumn{3}{c}{相邻月}\\口径 & 对数 & 满足条件 & 重合率 & 对数 & 满足条件 & 重合率\\\midrule
账务金额近稳 & 582 & 576 & 98.97\% & 439 & 438 & 99.77\%\\
还款与交易状态 & 1063 & 404 & 38.01\% & 943 & 325 & 34.46\%\\
支付与交易结构 & 403 & 271 & 67.25\% & 284 & 193 & 67.96\%\\
\bottomrule\end{tabularx}
\tabnote{满足条件且当期预测超出允许区间时才发生候选分修正；容量保护再检查入选边界。}
\end{table}
\fig{v21_business_tail}{不同业务状态下的尾部幅度与配对差。a展示本文、S、T；b为本文减对手，误差线采用42项比较族的同时95\%区间。负值表示尾部波动减小。}
两类业务客群与七字段门控的重合率为34.46\%--67.96\%。相对S的四项CVaR95差值，同时95\%区间均低于0。账务近稳组相邻月的逐项区间低于0，同时区间覆盖0。三类客群、两期结果和全部对手统一纳入比较族。

\clearpage\subsection*{A.9（续）账务近稳组的完整配对区间}
\begin{table}[H]\centering\caption{账务金额近稳的CVaR95配对差与区间}
\fontsize{9bp}{12bp}\selectfont
\begin{tabularx}{\linewidth}{@{}Yrrr@{}}\toprule
阶段 / 对手 & 本文减对手 & 逐项95\%区间 & 同时95\%区间\\\midrule
提额月 / 目标加权 R & -11.830 & [-14.506, -9.004] & [-17.166, -6.493]\\
提额月 / 风险集成 S & -8.810 & [-10.936, -6.558] & [-13.167, -4.454]\\
提额月 / 尾部优化 T & -1.450 & [-2.433, -0.475] & [-3.369, 0.468]\\
提额月 / 参考混合 FX & -1.794 & [-2.706, -0.701] & [-3.788, 0.199]\\
提额月 / 参考集成 E & -1.865 & [-2.949, -0.735] & [-4.091, 0.361]\\
提额月 / 同底模静态 & -5.194 & [-6.793, -3.533] & [-8.439, -1.950]\\
提额月 / 原1pp更新 & 0.852 & [0.097, 2.251] & [-1.698, 3.403]\\
相邻月 / 目标加权 R & -2.883 & [-5.481, -0.858] & [-7.591, 1.824]\\
相邻月 / 风险集成 S & -2.793 & [-5.338, -1.053] & [-7.204, 1.618]\\
相邻月 / 尾部优化 T & -0.030 & [-0.932, 0.860] & [-1.805, 1.746]\\
相邻月 / 参考混合 FX & -0.384 & [-1.253, 0.187] & [-1.896, 1.128]\\
相邻月 / 参考集成 E & -0.301 & [-1.132, 0.259] & [-1.755, 1.154]\\
相邻月 / 同底模静态 & -1.880 & [-4.576, -0.162] & [-6.735, 2.975]\\
相邻月 / 原1pp更新 & 0.041 & [0.000, 0.123] & [-0.107, 0.189]\\
\bottomrule\end{tabularx}
\tabnote{单位为概率百分点。2,000次客户共同重采样；同时区间覆盖三客群、两期、七对手共42项，临界值3.8839。}
\end{table}
CVaR95按经验分布最高5\%质量积分，边界样本按剩余质量计入。令$D_{(1)}\geq\cdots\geq D_{(n)}$，$m=0.05n$、$k=\lfloor m\rfloor$，则
\[
\widehat{\operatorname{CVaR}}_{.95}=\frac{\sum_{i=1}^{k}D_{(i)}+(m-k)D_{(k+1)}}{m}.
\]
P95定位波动分布的95\%分位，CVaR95衡量最高5\%尾部的平均幅度。总体分布变化由PSI监测，同一客户的跨期稳定性由配对变化衡量。

\clearpage\subsection*{A.9（续）还款状态与支付结构组}
\begin{table}[H]\centering\caption{还款与交易状态的CVaR95配对差与区间}
\fontsize{9bp}{12bp}\selectfont
\begin{tabularx}{\linewidth}{@{}Yrrr@{}}\toprule
阶段 / 对手 & 本文减对手 & 逐项95\%区间 & 同时95\%区间\\\midrule
提额月 / 目标加权 R & -6.150 & [-8.161, -4.213] & [-10.098, -2.202]\\
提额月 / 风险集成 S & -3.446 & [-5.045, -1.968] & [-6.454, -0.439]\\
提额月 / 尾部优化 T & -0.715 & [-1.488, 0.057] & [-2.230, 0.800]\\
提额月 / 参考混合 FX & -0.584 & [-1.032, -0.094] & [-1.520, 0.351]\\
提额月 / 参考集成 E & -0.611 & [-1.419, 0.158] & [-2.181, 0.959]\\
提额月 / 同底模静态 & -1.883 & [-2.862, -1.033] & [-3.697, -0.069]\\
提额月 / 原1pp更新 & 0.355 & [0.000, 1.074] & [-0.980, 1.690]\\
相邻月 / 目标加权 R & -1.879 & [-3.034, -0.744] & [-4.138, 0.381]\\
相邻月 / 风险集成 S & -2.992 & [-4.103, -1.889] & [-5.228, -0.757]\\
相邻月 / 尾部优化 T & -0.671 & [-1.152, -0.155] & [-1.654, 0.311]\\
相邻月 / 参考混合 FX & -0.125 & [-0.222, -0.022] & [-0.326, 0.076]\\
相邻月 / 参考集成 E & -0.022 & [-0.565, 0.504] & [-1.055, 1.011]\\
相邻月 / 同底模静态 & -0.412 & [-0.838, -0.101] & [-1.168, 0.343]\\
相邻月 / 原1pp更新 & 0.000 & [0.000, 0.001] & [-0.004, 0.004]\\
\bottomrule\end{tabularx}
\tabnote{单位为概率百分点。2,000次客户共同重采样；同时区间覆盖三客群、两期、七对手共42项，临界值3.8839。}
\end{table}
\begin{table}[H]\centering\caption{支付与交易结构的CVaR95配对差与区间}
\fontsize{9bp}{12bp}\selectfont
\begin{tabularx}{\linewidth}{@{}Yrrr@{}}\toprule
阶段 / 对手 & 本文减对手 & 逐项95\%区间 & 同时95\%区间\\\midrule
提额月 / 目标加权 R & -8.635 & [-11.951, -5.395] & [-15.149, -2.121]\\
提额月 / 风险集成 S & -6.897 & [-9.929, -4.003] & [-12.734, -1.060]\\
提额月 / 尾部优化 T & -1.667 & [-2.799, -0.583] & [-3.861, 0.526]\\
提额月 / 参考混合 FX & -1.558 & [-2.452, -0.729] & [-3.294, 0.178]\\
提额月 / 参考集成 E & -1.169 & [-2.426, -0.055] & [-3.545, 1.207]\\
提额月 / 同底模静态 & -4.092 & [-6.283, -2.105] & [-8.294, 0.109]\\
提额月 / 原1pp更新 & 0.048 & [0.000, 0.204] & [-0.169, 0.266]\\
相邻月 / 目标加权 R & -2.388 & [-4.344, -0.796] & [-5.835, 1.060]\\
相邻月 / 风险集成 S & -4.080 & [-5.795, -2.229] & [-7.650, -0.509]\\
相邻月 / 尾部优化 T & -0.607 & [-1.433, 0.136] & [-2.170, 0.956]\\
相邻月 / 参考混合 FX & 0.033 & [-0.185, 0.361] & [-0.545, 0.611]\\
相邻月 / 参考集成 E & -0.577 & [-1.778, 0.344] & [-2.664, 1.510]\\
相邻月 / 同底模静态 & -0.462 & [-0.929, -0.060] & [-1.329, 0.405]\\
相邻月 / 原1pp更新 & 0.000 & [-0.007, 0.077] & [-0.092, 0.092]\\
\bottomrule\end{tabularx}
\tabnote{单位为概率百分点。2,000次客户共同重采样；同时区间覆盖三客群、两期、七对手共42项，临界值3.8839。}
\end{table}

\clearpage\subsection*{A.10 容量规则选择与方向性比较}
\begin{table}[H]\centering\caption{开发集上的容量保护方式}
\fontsize{9bp}{12bp}\selectfont
\begin{tabularx}{\linewidth}{@{}Yrrrrr@{}}\toprule
保护方式 & 平均AUC & 平均AP & 10\%覆盖 & 较差期CVaR95 & 较差期P95\\\midrule
10\%边界 & 0.8004 & 0.1672 & 173 & 1.612 & 1.000\\
5\%边界 & 0.8006 & 0.1680 & 172 & 2.097 & 1.000\\
8\%边界 & 0.8004 & 0.1672 & 172 & 1.658 & 1.000\\
12\%边界 & 0.8003 & 0.1671 & 173 & 1.489 & 1.000\\
15\%边界 & 0.8004 & 0.1672 & 173 & 1.474 & 1.000\\
20\%边界 & 0.8004 & 0.1672 & 172 & 1.377 & 1.000\\
六档同时保护 & 0.8007 & 0.1682 & 173 & 2.929 & 1.214\\
\bottomrule\end{tabularx}
\tabnote{提额月6,411条、相邻月6,308条交叉拟合更新。各方案保留自身指定边界；主业务容量预定为10\%。幅度单位为概率百分点。}
\end{table}
主业务预定10\%复核容量。保护该边界后，开发两期覆盖173个事件窗口，与同底模静态一致，较差期CVaR95为1.612；六档容量同时保护的覆盖相同，较差期CVaR95为2.929。增加保护边界会收紧可用更新空间，机构可依据分级复核需求选择容量规则。

开发阶段另比较“风险上升自由更新、仅限制下降”的方向性方案，同样保护10\%名单。其两期平均AUC为0.8008、AP为0.1681，覆盖173例；较差期CVaR95为5.606，提额月P95为2.398。放开上升方向保留了更多细微概率响应，也扩大了尾部波动。综合两类目标，主方案采用近稳候选与容量保护的组合。

全排序投影比较采用更强约束：同时固定当前底模次序、非近稳客户读数及以上期静态分为中心的近稳区间。在这组条件下，1pp共同区间不可行，两期最小可行半径为21.94、25.57pp。容量保护因而保留业务必需的入选边界，为其余排序留出更新空间。

\subsection*{批次管理与跨期比较}
最终评分在同一模型、参考周期和容量版本下比较。扩大候选名单、调整容量或引入新账户时，系统整批重算并生成新批次版本。名单保持以指定底模和预定容量为准；AUC、AP、概率损失及完整排序均按最终$q$独立评价。

\clearpage\subsection*{A.11 输入适用范围与账务响应}
输入范围检查使用107,880个拟合行，其中一般窗口89,433行、事件开发18,447行。414项特征中，396项观测充分且可变。各维取训练1\%--99\%区间，每行以有值特征为分母计算越界比例；高越界阈值取训练行P95，即5.089\%。该指标衡量逐维偏离程度，联合支持结构另行估计。

\fig{v21_input_support}{参考尺度下的余额扩张与输入范围。a为余额超过参考额度的客户比例；b为特征越界比例超过训练P95的客户比例。工作台同时呈现原始余额、实际额度和当前使用率。}
\begin{table}[H]\centering\caption{参考尺度输入检查与分层风险识别}
\fontsize{9bp}{12bp}\selectfont
\begin{tabularx}{\linewidth}{@{}Yrrrrr@{}}\toprule
阶段 / 分组 & 人数 & 已知 & 事件 & 本文AUC & S的AUC\\\midrule
提额月 / 余额超过参考额度 & 937 & 919 & 55 & 0.7611 & 0.7622\\
提额月 / 高越界比例 & 179 & 172 & 5 & 0.9329 & 0.8958\\
提额月 / 其余输入 & 2728 & 2669 & 94 & 0.8096 & 0.8142\\
相邻月 / 余额超过参考额度 & 1268 & 1247 & 62 & 0.7503 & 0.7435\\
相邻月 / 高越界比例 & 268 & 264 & 7 & 0.7788 & 0.8071\\
相邻月 / 其余输入 & 2409 & 2359 & 81 & 0.7695 & 0.7723\\
\bottomrule\end{tabularx}
\tabnote{396个可变特征逐维1\%--99\%区间；高越界比例以训练行P95为阈值。“其余输入”为高越界组补集；余额超过参考额度可与后二组重合。}
\end{table}
提额月、相邻月分别有937、1,268名客户的余额超过参考额度；相邻月该组包含62／88个事件。参考比率大于1反映相对旧尺度的扩张，真实超限仍按当期额度认定。输入高越界记录交模型与数据岗位复核。

月内支付由覆盖最低应还转为未覆盖的记录有195／217对，其中188／211对结果已知，未来事件为2／1例；本文在该组的CVaR95为2.535／2.439。月度表未记录账单日、到期日和入账顺序，支付差额因此转入账务核对，补齐时间与凭证后形成办理意见。
\clearpage\subsection*{A.12 复核作业中心：查账、核证与责任接续}
作业中心围绕一件客户复核任务，汇集六期账务、参考额度和证据状态。A提供原始账务，B增加余额、额度、支付及最低应还趋势，C进一步提供参考尺度分解与带来源的材料草稿。经办人员切换记录、趋势和证据视图时，办理表持续保留已核读数、引用及下一岗位事项。图A7以同一练习案件展示三种工作方式。

\fig{trial_operations}{从账务核对到材料接续。a读取L1--L6原始记录；b对照余额、实际额度及还款轨迹；c分解余额与尺度带来的读数变化，并连接证据来源和待办事项。三种视图统一记录案件领取、提交、复审与返工。}

\clearpage\subsection*{A.12（续）质量、效率与接入成本的联合验收}
平台配置18个正式案例和3个独立练习案例。正式案例按“余额升而使用率降”“月内支付新差额”“其余情形”各选6个，均依据评分时可见信息取样。三种界面共享冻结案例、原始字段与评分版本；主工作台展示最新容量保护评分，作业中心比较不同界面对核对质量和材料接续的影响。

\textbf{以合格材料为效率比较的前提。}合格复核须做到读数正确、判断有来源、待办有责任。主比较C对B衡量参考尺度与证据协同的增量，A作为原始账务对照。每位参与者办理9案，每种界面3案，同人同案只出现一次。按岗位和经验分层轮换6种界面顺序；每个完整6人区组使同案在各界面各出现2次，下一区组换用另外9案。

\begin{table}[H]\centering\caption{作业验收指标与业务含义}\small\renewcommand{\arraystretch}{1.2}
\begin{tabularx}{\linewidth}{@{}p{24mm}Yp{29mm}@{}}\toprule
验收项目 & 判定和计算 & 对应业务目标\\\midrule
读数正确率 & 余额差、实际与参考使用率、支付关系、动作证据状态五项全部正确 & 关键读数一次核准\\\addlinespace[3bp]
材料完整正确率 & 读数及必填材料合格，证据完整性与处理适当性均获2分，最终复核通过 & 下一岗位直接接续\\\addlinespace[3bp]
合格材料工时 & 全部投入时间（含失败、返工）除以完整正确案件数；零合格件时单列 & 衡量完整处理成本\\\addlinespace[3bp]
返工与接续质量 & 退回原因、补件次数、重复取数、下一责任岗及来源定位 & 减少重复核对\\\bottomrule
\end{tabularx}\end{table}

\textbf{按完整办理过程计量成本。}正式案件在点击开始后展示，同时记录前台有效时间与墙钟时间。提交时保存版本，重复请求去重，刷新后恢复办理内容；退回后沿原案续办，累计全部投入时间。独立复审隐藏界面条件和参与者身份，按统一标准评价材料。

\textbf{让质量评价落到证据与处置。}证据完整性分0--2级：引用错误或缺失为0，主要判断有来源但缺关键期次为1，关键判断可定位至前后期、参考额度和凭证状态为2。处理适当性同分三级：将读数变化直接视为信用改善或直接执行调额为0，仅描述现象为1，区分读数与风险、明确待核凭证及下一责任岗为2。评分分歧经协商解决，初始意见同时保留。

\textbf{用现场结果决定试点扩围。}平台与案例已就绪，拟开展24名从业者先导试点。试点将回收正确率、工时和返工记录，报告配对差异，并按参与者成组重采样计算区间。首轮结论对应固定18案，扩围时按业务类型重新抽样。建议门槛为材料完整正确率差的95\%下限不低于$-5$个百分点，合格材料单位工时差的95\%上限小于0；业务负责人在试点前确认门槛，结果待现场回收。

\textbf{把采购成本与日常作业连接起来。}年化净作业价值按实际案件量、每案节时、全成本小时费率及年度总成本计算。成本包括接入摊销、维护、推理、监测、培训与人工复核。技术验收覆盖观察时点、字段映射、评分一致性、岗位权限、工单回执和恢复流程；风险验收回收动作原因、审批时间、后续账务与到期结局。228名连续入选客户体现了沿原案接续的现实需求，重复取数与材料重做的改善以现场记录计量。
